% theme: water_pressure_force_composition_decomposition
\MajorQuestion{水圧／力の合成・分解}
\SetRowSpaceQanda

\Instruction{水中ではたらく圧力に関する次の問いに答えなさい。}
\QQ{空気の重さによって生じる圧力をなんと呼ぶか。}{気圧}
\QQ{水の重さによって、水中の物体にはたらく圧力を\NamedBlank{water_pressure}という。}{水圧}
\QQ{\RefBlank{water_pressure}は、水中の物体にどの向きにはたらくか。}{あらゆる向き}
\QQ{水面からの深さが深くなるほど、\RefBlank{water_pressure}の大きさはどうなるか。}{大きくなる}
\QQ{\RefBlank{water_pressure}は、物体の表面に対してどの向きにはたらくか。}{垂直}

\Instruction{水中の物体にはたらく上向きの力に関する次の問いに答えなさい。}
\QQ{物体の上面にはたらく下向きの\RefBlank{water_pressure}と、下面にはたらく上向きの\RefBlank{water_pressure}の差によって\NamedBlank{buoyancy}が生じる。}{浮力}
\QQ{物体全体が水中に沈んでいるとき、水面からの深さが深くなると、\RefBlank{buoyancy}の大きさはどうなるか。}{変化しない}
\QQ{水中で物体全体が沈んでいるとき、物体の体積が大きいほど\RefBlank{buoyancy}はどうなるか。}{大きくなる}
\QQ{水中の物体にはたらく\RefBlank{buoyancy}が重力より大きいとき、物体はどうなるか。}{浮き上がる}
\QQ{浮力の単位は何か。}{N(ニュートン)}

\Instruction{複数の力を1つにまとめる操作に関する次の問いに答えなさい。}
\QQ{複数の力と同じはたらきをする1つの力を\NamedBlank{resultant_force}という。}{合力}
\QQ{複数の力から\RefBlank{resultant_force}を求めることを\NamedBlank{force_composition}という。}{力の合成}
\QQ{反対向きの2つの力を\RefBlank{force_composition}すると、\RefBlank{resultant_force}の大きさは2つの力の何になるか。}{差}
\QQ{\RefBlank{resultant_force}は、角度をもってはたらく2力がなす平行四辺形の\NamedBlank{taikaku}になる。}{対角線}
\QQ{\RefBlank{resultant_force}が、角度をもってはたらく2力がなす平行四辺形の\RefBlank{taikaku}になる法則をなんと言うか。}{力の平行四辺形の法則}

\Instruction{1つの力を分ける操作に関する次の問いに答えなさい。}
\QQ{1つの力を、それと同じはたらきをする2つの力に分けることを\NamedBlank{force_decomposition}という。}{力の分解}
\QQ{\RefBlank{force_decomposition}によってできた、それぞれの力を\NamedBlank{component_force}という。}{分力}
\QQ{もとの力の大きさが変わらない場合、2つの同じ大きさの\RefBlank{component_force}のなす角が小さいほど、\RefBlank{component_force}の大きさはどうなるか。}{小さくなる}
\QQ{斜面などがその上にある物体を支える力を\NamedBlank{kouryoku}と呼ぶ。}{垂直抗力}
\QQ{斜面が急になると、その上にある物体に働く\RefBlank{kouryoku}の大きさはどうなるか。}{小さくなる}
